\documentclass[11pt,a4paper]{article}
\usepackage[utf8]{inputenc}
\usepackage[T1]{fontenc}
\usepackage{lmodern}
\usepackage[french]{babel}
\usepackage{amsmath,amssymb,mathtools}
\usepackage{icomma}
\usepackage[version=4]{mhchem}
\usepackage{siunitx}
\sisetup{output-decimal-marker={,},per-mode=symbol}
\usepackage{xcolor}
\usepackage{colortbl}
\usepackage{tikz}
\usepackage[most]{tcolorbox}
\usepackage{enumitem}
\usepackage{array}
\usepackage{booktabs}
\usepackage{fancyhdr}
\usepackage[hidelinks]{hyperref}
\usetikzlibrary{arrows.meta,positioning,calc}
\usepackage[a4paper,margin=2.1cm,headheight=26pt,headsep=10pt]{geometry}

\definecolor{primary}{RGB}{146,95,161}
\definecolor{primarydark}{RGB}{92,55,108}
\definecolor{acidecol}{RGB}{205,60,45}
\definecolor{basecol}{RGB}{40,110,200}
\definecolor{excol}{RGB}{0,135,90}
\definecolor{attncol}{RGB}{200,40,55}

\tcbset{boxbase/.style={enhanced,breakable,boxrule=0.9pt,arc=2.5pt,left=3mm,right=3mm,top=2.3mm,bottom=2.3mm,fonttitle=\bfseries,coltitle=white,toptitle=0.6mm,bottomtitle=0.6mm}}
\newtcolorbox[auto counter]{corrbox}[1][]{boxbase,colback=excol!4,colframe=excol,title={\textsc{Corrigé}~\thetcbcounter\ifx\relax#1\relax\relax\else\textmd{ --- \itshape #1}\fi}}
\newtcolorbox{astuce}{boxbase,colback=attncol!6,colframe=attncol,sharp corners}

\pagestyle{fancy}
\fancyhf{}
\fancyhead[L]{\small\textcolor{primarydark}{\textbf{Fiche 10} — Thème 2 : pH, pOH, $K_e$}}
\fancyhead[R]{\small\textcolor{primarydark}{\textsc{Corrigé — Difficile}}}
\fancyfoot[C]{\small\thepage}
\renewcommand{\headrulewidth}{0.8pt}
\renewcommand{\headrule}{\hbox to\headwidth{\color{primary}\leaders\hrule height \headrulewidth\hfill}}

\setlength{\parindent}{0pt}
\setlength{\parskip}{3pt}

\begin{document}

\begin{center}
{\LARGE\bfseries\color{primarydark} Thème 2 --- Corrigé Difficile}
\end{center}

\begin{corrbox}[le pH neutre se déplace avec la température]
\begin{enumerate}[label=\alph*),leftmargin=1.6em,itemsep=2pt]
  \item L'autoprotolyse $\ce{2 H2O <=> H3O+ + OH-}$ est \textbf{endothermique} : chauffer favorise le sens de
        production d'ions (loi de Le Chatelier), donc $K_e$ augmente avec $T$.
  \item Milieu neutre : $[\ce{H3O+}]=[\ce{OH-}]$, donc $K_e=[\ce{H3O+}]^2$ et
        $[\ce{H3O+}]=\sqrt{10^{-13}} = 10^{-6{,}5}\approx 3{,}2\times10^{-7}\ \si{mol.L^{-1}}$.\\
        $\mathrm{pH}_{\text{neutre}} = -\log(10^{-6{,}5}) = \mathbf{6{,}5}$.
  \item À cette température, le neutre est à $6{,}5$ : une solution à $\mathrm{pH}=7$ ($>6{,}5$) est
        \textbf{basique} ; une solution à $\mathrm{pH}=6{,}5$ est exactement \textbf{neutre}.
  \item La frontière neutre est le pH de l'eau pure, qui vaut $-\log\sqrt{K_e}$ : elle ne vaut 7 qu'à
        \SI{25}{\celsius}. Dès que la température change, « basique » ne correspond plus à « $\mathrm{pH}>7$ »
        mais à « $\mathrm{pH}>\mathrm{pH}_{\text{neutre}}$ ».
\end{enumerate}
\end{corrbox}

\begin{corrbox}[du suc gastrique au sang]
\begin{enumerate}[label=\alph*),leftmargin=1.6em,itemsep=2pt]
  \item Suc gastrique : $[\ce{H3O+}] = 10^{-1{,}5} \approx 3{,}2\times10^{-2}\ \si{mol.L^{-1}}$.\\
        Sang : $[\ce{H3O+}] = 10^{-7{,}4} \approx 4{,}0\times10^{-8}\ \si{mol.L^{-1}}$.
  \item Facteur $= 10^{\,7{,}4-1{,}5} = 10^{5{,}9} \approx \mathbf{8\times10^{5}}$ : le suc gastrique est presque
        un million de fois plus concentré en \ce{H3O+} que le sang.
  \item $[\ce{OH-}]_{\text{sang}} = \dfrac{K_e}{[\ce{H3O+}]} = \dfrac{10^{-14}}{4{,}0\times10^{-8}}
        \approx 2{,}5\times10^{-7}\ \si{mol.L^{-1}}$.\\
        Cohérence : $\mathrm{pOH} = 14 - 7{,}4 = 6{,}6$, et $10^{-6{,}6}\approx 2{,}5\times10^{-7}$. \checkmark\
        Le sang est légèrement basique ($[\ce{OH-}]>[\ce{H3O+}]$).
\end{enumerate}
\end{corrbox}

\begin{corrbox}[{démontrer le critère « acide $\Rightarrow [\ce{H3O+}]>[\ce{OH-}]$ »}]
\begin{enumerate}[label=\alph*),leftmargin=1.6em,itemsep=2pt]
  \item « Acide » $\Leftrightarrow \mathrm{pH}<7 \Leftrightarrow -\log[\ce{H3O+}]<7 \Leftrightarrow
        [\ce{H3O+}] > 10^{-7}\ \si{mol.L^{-1}}$.
  \item De $[\ce{OH-}] = \dfrac{K_e}{[\ce{H3O+}]} = \dfrac{10^{-14}}{[\ce{H3O+}]}$ et $[\ce{H3O+}]>10^{-7}$,
        on tire $[\ce{OH-}] < \dfrac{10^{-14}}{10^{-7}} = 10^{-7}\ \si{mol.L^{-1}}$.
  \item Donc $[\ce{H3O+}] > 10^{-7} > [\ce{OH-}]$, c'est-à-dire \textbf{$[\ce{H3O+}]>[\ce{OH-}]$}. Le
        raisonnement reste valable à toute température, à condition de remplacer le nombre $10^{-7}$ par
        $\sqrt{K_e}$ (la valeur « 7 » est propre à \SI{25}{\celsius}).
\end{enumerate}
\end{corrbox}

\begin{astuce}
\textbf{Ce que le concours attend.} Ne jamais réciter « acide = pH$<$7 » comme un dogme : savoir le
\emph{redémontrer} via $K_e$, et savoir que la borne dépend de la température. C'est la nuance qui distingue
les bonnes copies (cf.\ annales sur la solution « basique » à pH 7).
\end{astuce}

\end{document}
